% !TeX root = ../../../main.tex
\subsection{Imputation Methods Details} \label{ch04:app:imputation}

We follow \citet{qi2023an} to implement three surrogate imputation strategies for estimating the true event time \( T \) in right-censored survival data. Let \( Y = \min(T, C) \) be the observed time, with censoring indicator \( \delta = \mathbbm{1}\{T \leq C\} \). Let \( S_{\text{KM}(\mathcal{D})}(t) \) denote the Kaplan-Meier estimate of the survival function using the dataset \( \mathcal{D} \), and \( N \) the number of subjects. The three methods below are used to impute a surrogate outcome \( \tilde{T}_i \) for censored subject \( i \) observed at time \( t_i \).

\ChFourParagraph{1. Margin imputation:}
This method assigns a “best guess” value to each censored subject using the nonparametric Kaplan-Meier estimator. This surrogate value, called the \emph{margin time}, can be interpreted as the conditional expectation of the event time given that the event occurs after the censoring time. For a subject censored at time \( t_i \), the margin-imputed event time is computed as:
\begin{equation}
\tilde{T}^{\text{margin}}_i 
= \mathbb{E}[T_i \mid T_i > t_i] 
= t_i + \frac{\int_{t_i}^{\infty} S_{\text{KM}(\mathcal{D})}(t) \, dt}{S_{\text{KM}(\mathcal{D})}(t_i)}
\end{equation}
where \( S_{\text{KM}(\mathcal{D})}(t) \) is the Kaplan-Meier survival estimate derived from the training dataset.

The reliability of this imputation depends on the censoring time. For example, if a subject is censored very early (e.g., at time 0), the margin time is highly uncertain due to the lack of observed data beyond that point. In contrast, if a subject is censored near the maximum observed follow-up, the margin time is more likely to be close to the true event time.

\ChFourParagraph{2. IPCW-T imputation:}
This method imputes a surrogate event time for censored subjects based on the observed outcomes of subsequent uncensored individuals. Specifically, for a subject censored at time \( t_i \), the imputed value is calculated as the average event time of all uncensored subjects with observed times after \( t_i \):
\begin{equation}
\tilde{T}^{\text{IPCW}}_i = 
\frac{\sum_{j=1}^N \mathbbm{1}\{t_i < t_j\} \cdot \mathbbm{1}\{\delta_j = 1\} \cdot t_j}
     {\sum_{j=1}^N \mathbbm{1}\{t_i < t_j\} \cdot \mathbbm{1}\{\delta_j = 1\}}
\end{equation}
This imputes the event time for subject \( i \) by averaging the observed event times of those uncensored subjects who experienced the event after \( t_i \). The method is motivated by the idea that these subsequent subjects provide empirical evidence about the possible timing of the unobserved event.

However, a limitation of this approach is that it fails to provide an imputation when there are no uncensored subjects observed after \( t_i \). In such cases, the denominator of the expression becomes zero, and the method is unable to approximate the event time. In \citet{qi2023an}, subjects for whom this occurs are excluded from evaluation, whereas in our setup we used the original observed time as the imputed time.


\ChFourParagraph{3. Pseudo-observation imputation:}
This method imputes the event time using pseudo-observations, which estimate the contribution of each subject to an overall unbiased estimator of the event time distribution. Let \( \hat{\theta} \) be an estimator of the mean event time based on right-censored data, and let \( \hat{\theta}^{-i} \) denote the same estimator computed with the \( i \)-th subject removed from the dataset. Then, the pseudo-observation for subject \( i \) is defined as:
\begin{equation}
\tilde{T}^{\text{pseudo}}_i = e_{\text{Pseudo-Obs}}(t_i, \mathcal{D}) 
= N \cdot \hat{\theta} - (N - 1) \cdot \hat{\theta}^{-i}
\end{equation}
This quantity can be interpreted as the individual contribution of subject \( i \) to the overall estimate \( \hat{\theta} \). In practice, both \( \hat{\theta} \) and \( \hat{\theta}^{-i} \) can be computed using the mean of the Kaplan-Meier survival curve:
\[
\hat{\theta} = \mathbb{E}_t[S_{\text{KM}(\mathcal{D})}(t)], \quad
\hat{\theta}^{-i} = \mathbb{E}_t[S_{\text{KM}(\mathcal{D} \setminus \{i\})}(t)]
\]
Once the pseudo-observations \( \tilde{T}^{\text{pseudo}}_i \) are computed for all censored subjects, they are substituted in place of the true event times for evaluation or modeling. 

Although pseudo-observations are not exact conditional expectations, they can approximate \( \mathbb{E}[T_i \mid X_i] \) under certain assumptions. In particular, when censoring is independent of covariates and the sample size is large, pseudo-observations behave asymptotically like i.i.d. draws from the true conditional expectation:
\[
\mathbb{E}[\tilde{T}^{\text{pseudo}}_i \mid X_i] \approx \mathbb{E}[T_i \mid X_i]
\]
This makes the pseudo-observation method a principled, nonparametric approach for imputing censored survival times, particularly when estimating global quantities like the mean event time.

These imputation strategies enable us to transform the survival outcome into a fully observed target variable, allowing the application of standard regression-based methods in causal effect estimation. 
To ensure meaningful estimates, it is important that each imputed event time for a censored subject is guaranteed to be greater than or equal to the censoring time---reflecting the fact that the true event must occur after the last time it was observed. 
In our implementation, we manually enforce this constraint by setting the imputed value to the observed censoring time whenever the imputation procedure yields a value less than \( t_i \).

%%%%%%%%%%%%%%%%%%%%%%%%%%%%%%%%%%%%%%%%%%%%%%%%%%%%%%%%%%%%
%%%%%%%%%%%%%%%%%%%%%%%%%%%%%%%%%%%%%%%%%%%%%%%%%%%%%%%%%%%%
%%%%%%%%%%%%%%%%%%%%%%%%%%%%%%%%%%%%%%%%%%%%%%%%%%%%%%%%%%%%



